% problem-000411 4 汞量法测定胆离子
\section*{4. 汞量法测定胆离子}
\textit{下列为分问。}

\subsection*{4-1}
简述配制硝酸汞溶液时必须⽤硝酸酸化的理由。

\subsection*{4-2}
称取 1.713 g $\mathrm { H g ( N O _ { 3 } ) _ { 2 } } { \cdot } x \mathrm { H _ { 2 } O }$ ，配制成 500 mL 溶液作为滴定剂。取 20.00 mL 0.0100 mol $\mathrm { L } ^ { \mathrm { - 1 } } \mathrm { N a C l }$ 标准溶液注⼊锥形瓶，⽤1 mL 5% $\mathrm { H N O _ { 3 } }$ 酸化，加⼊5滴⼆苯卡巴腙指⽰剂，⽤上述硝酸汞溶液滴定⾄紫⾊，消耗 $1 0 . 2 0 \mathrm { m L } _ { \circ }$ 推断该硝酸汞⽔合物样品的化学式。

\subsection*{4-3}
取0.500 mL⾎清放⼊⼩锥形瓶，加2mL去离⼦⽔、4滴5%的硝酸和3滴⼆苯卡巴腙指⽰剂，⽤上述硝酸汞溶液滴定⾄终点，消耗 $1 . 5 3 \mathrm { m L } _ { \circ }$ 为使测量结果准确，以⼗倍于⾎清样品体积的⽔为试样进⾏空⽩实验，消耗硝酸汞溶液0.80 mL。计算该⾎清样品中氯离⼦的浓度（毫克/100毫升）。
